\documentclass[12pt,a4paper]{article}
\usepackage[utf8]{inputenc}
\usepackage{xepersian}
\usepackage{graphicx}
\usepackage{hyperref}
\usepackage{listings}

\settextfont{XB Niloofar}
\setdigitfont{XB Niloofar}

\title{جریان درخواست‌ها و منابع ترافیک}
\author{}
\date{}

\begin{document}

\maketitle

\section{جریان درخواست‌ها}

\begin{enumerate}
\item کاربران از طریق نام‌دامنه‌ها مانند \texttt{api.mysite.com} به وب‌سایت‌ها دسترسی پیدا می‌کنند. معمولاً سیستم نام‌دامنه\footnote{DNS - Domain Name System} یک سرویس پولی است که توسط ارائه‌دهندگان شخص ثالث ارائه می‌شود و روی سرورهای ما میزبانی نمی‌شود.

\item آدرس پروتکل اینترنت\footnote{IP - Internet Protocol} به مرورگر یا اپلیکیشن موبایل بازگردانده می‌شود. در مثال ما، آدرس IP معادل ۱۵.۱۲۵.۲۳.۲۱۴ برگشت داده می‌شود.

\item پس از دریافت آدرس IP، درخواست‌های پروتکل انتقال ابرمتن\footnote{HTTP - Hypertext Transfer Protocol} مستقیماً به وب‌سرور شما ارسال می‌شوند.

\item وب‌سرور صفحات HTML یا پاسخ‌های JSON را برای نمایش بازمی‌گرداند.
\end{enumerate}

\section{منابع ترافیک}

ترافیک ورودی به وب‌سرور شما از دو منبع اصلی تأمین می‌شود:

\subsection{وب اپلیکیشن}
از ترکیبی از زبان‌های سمت سرور (مانند جاوا، پایتون و غیره) برای پردازش منطق کسب‌وکار، ذخیره‌سازی و... استفاده می‌کند و زبان‌های سمت کلاینت (HTML و جاوااسکریپت) برای نمایش محتوا به کار می‌روند.

\subsection{اپلیکیشن موبایل}
پروتکل HTTP، پروتکل ارتباطی بین اپلیکیشن موبایل و وب‌سرور است. فرمت رایج پاسخ API، نشانه‌گذاری شیء جاوااسکریپت\footnote{JSON - JavaScript Object Notation} است که به دلیل سادگی برای انتقال داده استفاده می‌شود.

مثال زیر یک پاسخ API با فرمت JSON را نشان می‌دهد:

\begin{latin}
\begin{lstlisting}
GET /users/12 – دریافت شیء کاربر مربوط به شناسه ۱۲
\end{lstlisting}
\end{latin}

\end{document}
